% !TEX root = ../main.tex
% =====================================================================
% Section 3 — Method  (gdn2vessel, ACCV 2026)
% Standalone draft. Intended to be \input into main.tex.
% Defensive-writing rules R1--R7 enforced (see STORY_FRAMEWORK.md).
%   R1: no "universal"/"always"; reconnection degrades with severity/length.
%   R2: no "we prove"/"theorem"/"uniqueness"; "we design / motivated by".
%   R3: GDKVM distinction handled in related_work.tex (kept consistent here).
%   R4: scan/reorder is a standard engineering choice, NOT a contribution.
%   R5: layered supervision honesty — backbone/memory/Frangi gate are
%       input-derived (never GT topology); re-ID head is synthetic-break weak
%       supervision, stop-gradient isolated, never relabelled as unsupervised.
%   R6: connectivity/re-ID numbers carry statistics (placeholders here).
%   R7: CV-method contribution first; medical datasets are validation.
% All performance numbers are placeholders (\TODO) pending verifier csv checks.
% =====================================================================

\section{Method}
\label{sec:method}

We design a single end-to-end segmentation model that augments a convolutional
U-Net with a delta-rule associative memory placed in the deep encoder, and reads
that memory out to re-identify the two sides of a vessel break. We first give the
overall architecture (Section~\ref{sec:method:arch}), then describe how the GDN-2
matrix state acts as a vessel-identity memory that supports in-model spatial
reconnection (Section~\ref{sec:method:memory}, Claim~1), the explicit spatial
re-identification mechanism and its supervision (Section~\ref{sec:method:reid},
Claim~2), the differentiable Frangi decoupled gate that modulates the memory
(Section~\ref{sec:method:frangi}, Claim~3), and finally the standard 2D scan that
turns the spatial feature map into the sequence consumed by the memory
(Section~\ref{sec:method:scan}).

% ---------------------------------------------------------------------
\subsection{Overall Architecture}
\label{sec:method:arch}

The backbone is a standard convolutional U-Net encoder--decoder. At a deep encoder
stage, where the spatial resolution is low and the receptive field already spans
local vessel context, we insert a GDN-2 associative-memory module. Let
$F \in \mathbb{R}^{C \times H \times W}$ be the bottleneck feature map. We flatten
$F$ into a token sequence $(x_1, \dots, x_L)$ with $L = H \cdot W$ using a standard
scan order (Section~\ref{sec:method:scan}); we keep $L \lesssim 1\mathrm{K}$ by
placing the module at a sufficiently downsampled stage. The memory module processes
this sequence, producing a memory-augmented feature map that is folded back into the
encoder tensor and passed to the decoder, so the whole model remains a single
forward pass.

In parallel, a lightweight differentiable Frangi layer consumes \emph{only} the
input image and produces a vesselness map that is downsampled to the bottleneck
resolution. This map is the external control signal for the memory's erase/write
gates (Section~\ref{sec:method:frangi}); it is never used as a label or as a
segmentation output. Finally, a re-identification read-out head takes pairs of
memory states sampled at candidate break endpoints and predicts whether they belong
to the same vessel (Section~\ref{sec:method:reid}). The segmentation decoder and the
re-identification head share the encoder and the memory features, but, as detailed
below, the re-identification supervision is stop-gradient isolated so that it updates
only the matching head and does not flow back into the memory or the encoder.

% ---------------------------------------------------------------------
\subsection{GDN-2 Associative Memory as Vessel-Identity Memory}
\label{sec:method:memory}

\paragraph{Mechanism.}
We use the recurrent matrix state of GDN-2 as a vessel-identity memory. As the
module scans the flattened bottleneck sequence, it writes per-token key--value
pairs into a matrix state $S_t$ and updates it with the gated delta rule. Writing
the per-token query, key and value as $q_t, k_t, v_t$ (obtained from $x_t$ via
learned projections), the state update follows the gated delta-rule recurrence
\begin{equation}
  S_t = S_{t-1}\,\mathrm{diag}(g_t)\;-\;\beta_t\,\bigl(S_{t-1} k_t - v_t\bigr)\,k_t^{\top},
  \qquad
  o_t = S_t\, q_t ,
  \label{eq:gdn2}
\end{equation}
where $g_t$ is a (channel-wise) decay/erase signal that controls how much of the
existing memory is retained, and $\beta_t$ is a write strength that controls how
strongly the current key--value pair overwrites the state. The output $o_t$
retrieves, for the current query, the value previously associated with similar keys.
Intuitively, while the scan traverses a vessel, the module writes that vessel's
appearance into $S_t$ under a key that encodes its local geometry; when the scan
reaches the far side of a low-contrast or occluded break, a query with a similar key
retrieves the earlier value, continuing the same structure in feature space rather
than starting a new one.

This is an \emph{in-model, end-to-end} reconnection mechanism. It differs from a
topological loss such as clDice (a global regulariser applied to the prediction),
from graph-based post-processing, and from learned post-processing that operates on
an already produced mask~\cite{creatis}: here the reconnection is realised inside the
forward pass through the evolving memory state, and its parameters are trained by the
segmentation objective end to end.

\paragraph{Scope of the claim (R1).}
We do \emph{not} claim universal or guaranteed reconnection. The memory has finite
capacity and the retrieval signal weakens as the break becomes more severe (wider
gap, thinner vessel) and as the number of intervening tokens between the two sides of
the break grows along the scan. We therefore characterise reconnection as a
capability that \emph{degrades with break severity and with sequence length}, and we
report this degradation explicitly as severity and length curves in our experiments
(Section~\TODO{ref to ablation}); we do not present a single headline number as if
reconnection were unconditional. The memory keys are derived purely from the input
features and the Frangi gate (Section~\ref{sec:method:frangi}); they never read
ground-truth topology, which keeps the Claim~1 reconnection ability
ground-truth-topology-free.

% ---------------------------------------------------------------------
\subsection{Spatial Re-Identification Mechanism}
\label{sec:method:reid}

Existing reconnection metrics, such as the normalised connected-component error
$\varepsilon_{\beta_0}$~\cite{creatis}, quantify only whether a gap has been
\emph{filled}; they do not test whether the filler is the continuation of the
\emph{same} vessel. We make this distinction explicit with a re-identification
read-out head and a corresponding re-identification rate metric, defined and
evaluated in Section~\ref{sec:bench}. Motivated by the identity-aware association
logic of multi-object tracking (e.g. IDF1), and following prior use of such
tracking-style measures for vessel tracing~\cite{snaketracker}, we score a
reconnection as correct only when the two sides of a break are matched to the same
underlying vessel, not merely connected.

\paragraph{Re-identification head.}
Given two candidate endpoints across a break, we sample their memory-augmented
features and feed the pair to a small matching head that predicts a same-vessel
score. At inference, the head decides, for each gap, whether the bridging belongs to
the same vessel, which yields the re-identification rate.

\paragraph{Supervision scope (R5, honest layering).}
The matching head is trained with \emph{synthetic-break weak supervision}. Because we
generate the breaks ourselves with a controlled procedure
(Section~\ref{sec:bench}), we know by construction which two sides originate from the
same vessel, and we use this as the same-vessel label for the head. This is the same
plug-and-play regime as the learned post-processing of Corvo~\emph{et
al.}~\cite{creatis}, which likewise does not claim to be annotation-free; we make no
unsupervised claim for the matching head. Critically, this weak supervision is
\emph{stop-gradient isolated}: the memory-augmented features fed to the head are
detached, so the same-vessel loss updates \emph{only} the matching head and does not
back-propagate into the memory or the encoder. Consequently the reconnection ability
of Claim~1 remains driven by the segmentation objective and the input-derived signals
alone, and is not learned from ground-truth topology. The novelty here is the
\emph{mechanism} (an associative memory repurposed for spatial re-identification) and
the \emph{benchmark}, not annotation-free learning.

\paragraph{Attribution design (pre-registered).}
To attribute the re-identification gain to the \emph{stateful} delta-rule memory
rather than to the extra capacity of adding an attention block, we pre-register a
three-arm comparison, all arms sharing the same seeds ($\geq 3$), splits,
hyper-parameters, and training budget, and differing only in the bottleneck feature
source feeding the same re-identification head:
\begin{itemize}
  \item \textbf{A0$'$ (null arm).} A plain CNN bottleneck with no memory and no
        attention module (\texttt{use\_memory=False}); the floor with neither
        capacity nor mechanism.
  \item \textbf{A1$'$ (iso-parameter arm).} A parameter-matched plain linear-attention
        module that keeps the same query/key/value/output projections,
        LayerNorm, insertion depth, and Frangi-gate pathway, but replaces the
        stateful delta-rule memory with a \emph{stateless} linear attention; the total
        number of trainable parameters is matched to A2 within $\pm 5\%$ (the measured
        counts are reported in a table footnote). This isolates the contribution of the
        added attention capacity.
  \item \textbf{A2 (full).} The GDN-2 stateful delta-rule associative memory, with the
        delta update and the recurrent state $S_t$; this is our full model.
\end{itemize}
We further include A3, which detaches the memory during training to separate the
contribution of the head from that of training the memory, and a leakage-control arm
A4 (in Section~\ref{sec:bench}) that supervises the head with \emph{predicted} rather
than ground-truth skeleton breaks. The headline attribution is the ordering
$\text{A2} > \text{A1}' > \text{A0}'$ in re-identification rate, with paired
significance reported per dataset; if A2 does not exceed the iso-parameter A1$'$, the
gain would be attributable to capacity rather than to the stateful associative memory,
and we report this honestly.

\paragraph{Re-identification beyond filling: a controlled direct effect.}
A natural concern is that a higher re-identification rate is merely a by-product of
better gap filling (lower $\varepsilon_{\beta_0}$). We treat $\varepsilon_{\beta_0}$
as a \emph{mediator} on the causal path
$\text{memory} \rightarrow \text{better filling} \rightarrow \text{re-identification}$,
not as a confounder; conditioning on a mediator would induce over-control bias, so we
do \emph{not} regress re-identification on memory while controlling for
$\varepsilon_{\beta_0}$. Instead we estimate a \emph{controlled direct effect} on a
filling-matched subset: because A2 and A1$'$ are naturally paired on the same image, we
retain only the images for which the two arms achieve comparable filling quality
(those whose $|\varepsilon_{\beta_0}(\mathrm{A2}) -
\varepsilon_{\beta_0}(\mathrm{A1}')|$ falls within the per-dataset
$\varepsilon_{\beta_0}$ inter-quartile range) and re-run the paired comparison within
that subset. If A2 still re-identifies more accurately when filling quality is matched,
the re-identification gain is a direct effect of the memory rather than a by-product of
$\varepsilon_{\beta_0}$. We additionally report an exploratory mediation decomposition
(total, natural direct, and natural indirect effects) for completeness; this
decomposition is descriptive only and is not used as a pass/fail criterion, since the
natural direct effect is expected to be attenuated when the mediating path is cut.
Statistical tests use exact paired sign/permutation procedures and bootstrap
confidence intervals implemented by hand (R6); we avoid relying on heavy library
defaults.

% ---------------------------------------------------------------------
\subsection{Differentiable Frangi Decoupled Gate}
\label{sec:method:frangi}

The GDN-2 memory exposes two control signals in Eq.~\eqref{eq:gdn2}: the decay/erase
signal $g_t$ and the write strength $\beta_t$. We design a \emph{decoupled} gate that
modulates these from an input-derived vesselness cue, so that along a vessel the model
is encouraged to retain its identity memory (suppress erase) and at fresh vessel
evidence it is encouraged to write (raise write). The modulating signal is the output
of a differentiable Frangi layer that consumes only the input image; it is never the
segmentation prediction or the ground truth, which keeps the gate input-derived and
avoids a chicken-and-egg dependence on the very segmentation we are producing.

\paragraph{Kernel-external modulation (R3/R4 honesty).}
The general gated delta-rule kernel we build on exposes a single write strength
$\beta$. We therefore realise the decoupled erase/write behaviour as a
\emph{kernel-external} modulation layer: the Frangi gate produces a write correction
that raises $\beta_t$ at vessel evidence and an additive decay correction that lowers
the erase of $g_t$ along a continuing vessel. We are explicit that this is a
\emph{two-gate approximation} layered on top of the single-$\beta$ kernel; we do
\emph{not} rewrite the kernel to implement two strictly independent gates, and we do
not claim exact equivalence to a true decoupled-kernel implementation. The
contribution is the associative-memory mechanism and its input-derived vesselness
modulation, not low-level kernel/scan engineering.

% ---------------------------------------------------------------------
\subsection{2D Scan and Multi-Directional Merge}
\label{sec:method:scan}

To apply a fundamentally sequential memory to a 2D feature map, we flatten the
bottleneck with a standard reorder and, as is common in vision SSM/linear-attention
backbones, run the memory over a small set of complementary scan directions and merge
their outputs. We adopt a standard scan-and-merge implementation; the contribution of
this paper is the associative memory and the spatial re-identification it enables, not
the scan ordering (R4). The number of directions is a design choice we ablate rather
than promote, and we make no claim that the particular reorder is novel.
